% !TEX program = xelatex
\documentclass{英语试卷}
\begin{document}
\examtitle{武汉市武昌区 2025 届高三年级上学期期末质量检测}{英语}
\par 武汉市武昌区 2025 届高三年级上学期期末质量检测\par
\begin{notice}
\par 1. 答卷前，考生务必将自己的姓名、准考证号填写在答题卡上。\par
\par 2. 回答选择题时，选出每小题答案后，用铅笔把答题卡上对应题目的答案标号涂黑。如需改动，用橡皮擦干净后，再选涂其他答案标号。回答非选择题时，将答案写在答题卡上，写在本试卷上无效。\par
\par 3. 考试结束后，将本试卷和答题卡一并交回，\par
\end{notice}
\parttitle{第一部分 听力（共两节，满分 30 分）}
\par 做题时，先将答案标在试卷上。录音内容结束后，你将有两分钟的时间将试卷上的答案转涂到答题卡上。\par
\sectitle{第一节（共5小题；每小题1.5分，满分7.5分）}
\par 听下面 5 段对话。每段对话后有一个小题，从题中所给的 A、B、C 三个选项中选出最佳选项，并标在试卷的相应位置。听完每段对话后，你都有 10 秒钟的时间来回答有关小题和阅读下一小题。每段对话仅读一遍。\par
\question{例：How much is the shirt?}
\begin{choices}(3)
  \item £ 19.15.
  \item £ 9.18.
  \item £ 9.15. 答案是 C,
\end{choices}
\question{1. Where will the man go later?}
\begin{choices}(3)
  \item To the gym.
  \item To the mall.
  \item To the school.
\end{choices}
\question{2. Why does the woman want to leave?}
\begin{choices}(3)
  \item She's in a rush.
  \item She's hungry.
  \item She's angry.
\end{choices}
\question{3. What is the woman's problem?}
\begin{choices}(3)
  \item She lost her train ticket.
  \item She got on the wrong train.
  \item She forgot her seat number.
\end{choices}
\question{4. How much will the man spend on the gift?}
\begin{choices}(3)
  \item \$20.
  \item \$30.
  \item \$50.
\end{choices}
\question{5. What is the conversation mainly about?}
\begin{choices}(3)
  \item The homework.
  \item A math class.
  \item A football game.
\end{choices}
\sectitle{第二节（共 15 小题；每小题 1.5 分，满分 22.5 分）}
\par 听下面 5 段对话或独白。每段对话或独白后有几个小题，从题中所给的 A、B、C 三个选项中选出最佳选项，并标在试卷的相应位置。听每段对话或独白前，你将有时间阅读各个小题，每小题 5 秒钟；听完后，各小题将给出 5 秒钟的作答时间。每段对话或独白读两遍。\par
\par 听第 6 段材料，回答第 6 至 7 题。\par
\question{6. What color is the man's car?}
\begin{choices}(3)
  \item Red.
  \item Blue.
  \item Yellow.
\end{choices}
\question{7. How long did the man's journey take?}
\begin{choices}(3)
  \item A quarter of an hour.
  \item An hour and a quarter.
  \item Three quarters of an hour.
\end{choices}
\par 听第 7 段材料，回答第 8 至 10 题。\par
\question{8. What's the man worried about?}
\begin{choices}(3)
  \item The cat is acting unusually.
  \item Sarah has been gone for a few days.
  \item The woman hasn't eaten anything yet.
\end{choices}
\question{9. Where is Sarah?}
\begin{choices}(3)
  \item At school.
  \item At a camp.
  \item At a hospital.
\end{choices}
\question{10. What is the probable relationship between the speakers?}
\begin{choices}(3)
  \item Father and daughter.
  \item Patient and doctor.
  \item Husband and wife.
\end{choices}
\par 听第 8 段材料, 回答第 11 至 13 题。\par
\question{11. How long did the woman teach at the village school?}
\begin{choices}(3)
  \item For one year.
  \item For two years.
  \item For three years.
\end{choices}
\question{12. Where is the woman from?}
\begin{choices}(3)
  \item A city in China.
  \item A village in China.
  \item The countryside of Italy.
\end{choices}
\question{13. What does the man think of the woman's stories?}
\begin{choices}(3)
  \item They're not realistic.
  \item They're we l-written.
  \item They're too emotional.
\end{choices}
\par 听第 9 段材料, 回答第 14 至 16 题,\par
\question{14. What is Thomas hoping to do before starting work?}
\begin{choices}(3)
  \item Travel the world.
  \item Go to university.
  \item Join the army.
\end{choices}
\question{15. What is Thomas probably interested in?}
\begin{choices}(3)
  \item Landforms.
  \item Countries.
  \item Satellites.
\end{choices}
\question{16. How many desks will Thomas visit?}
\begin{choices}(3)
  \item One.
  \item Two.
  \item Three.
\end{choices}
\par 听第 10 段材料, 回答第 17 至 20 题,\par
\question{17. What is the benefit of spreading rock?}
\begin{choices}(3)
  \item Air quality becomes better.
  \item Plants obtain more vitamins.
  \item Farmers can use their fields again.
\end{choices}
\question{18. Why is spreading rock better than planting trees?}
\begin{choices}(2)
  \item It is cheaper. 3. It is quicker.
  \item It is more useful.
\end{choices}
\question{19. Who could earn money from this method?}
\begin{choices}(3)
  \item The mineral mines.
  \item The government.
  \item The farmers.
\end{choices}
\question{20. How can the mines easily adapt to this change?}
\begin{choices}(3)
  \item By finding new areas to dig.
  \item By looking for different minerals.
  \item By changing where their waste goes.
\end{choices}
\parttitle{第二部分 阅读理解（共两节，满分 50 分）}
\sectitle{第一节（共15小题；每小题2.5分，满分37.5分）}
\par 阅读下列短文，从每题所给的 A、B、C 和 D 四个选项中，选出最佳答案。\par
\passagelabel{A}
\par Film "Rear Window" (Aug.23)\par
\par Seventy years after its original release, Alfred Hitchcock's thriller about a New York photojournalist who believes he has witnessed a murder gets a limited theatrical run. Cux critic praised the film as a moving portrait of being trapped, writing that "few films better capture the mental state that results from being stuck in a room for an extended period of time."\par
\par TV "Pachinko" (Aug.23)\par
\par Based on Min Jin Lee's historical novel. this drama series follows four generations of a Korean family who immigrate to Japan through different timelines. The second season begins with family leader Kim Sunja trying to lead her family through the final months of World War II. Our critic described the show as a "beautiful and slow-moving journey" through history.\par
\par Art "What It Becomes" (Aug.24)\par
\par The Whitney Museum of American Art in New York is highlighting new and rarely exhibited work from 11 artists — including David Hammons, Catherine Opie and Wendy RedStar — this show looks at the ways contemporary creators are expanding the boundaries of drawing, in content, medium and technique.\par
\par Videogame "Black Myth: Wukong" (Aug.20)\par
\par This action RPG, inspired by the 16fh-century Chinese novel "Journey to the West," puts players in the role of the book’s Monkey King, here called the Destined One, as they use a powerful staff to battle their way through Chinese mythology (神话). With stunning graphics and exciting gameplay, the game immerses players in a world full of legendary creatures and fierce enemies.\par
\question{21. What is the highlight of "Rear Window"?}
\begin{choices}(2)
  \item The originality of a suspense thriller.
  \item The portrait of an important character.
  \item The life experience of a photojournalist.
  \item The account of a psychological condition.
\end{choices}
\question{22. What does "Pachinko" and "Black Myth: Wukong" have in common?}
\begin{choices}(2)
  \item Both explore historical events.
  \item Both are adaptations of literary works.
  \item Both are praised for their beautiful visuals.
  \item Both include elements of fantasy and mythology
\end{choices}
\question{23. What type of writing is this passage?}
\begin{choices}(2)
  \item A review.
  \item A report.
  \item An editorial.
  \item An advertisement.
\end{choices}
\passagelabel{B}
\par In June 2024, Caitlin Clark, America's beloved college basketball sensation, faced a harsh reality check in her WNBA debut(首秀). Struggling with her shot, her handle, and adapting to a new team, Clark appeared overmatched, unable to bring forth the magic that helped the distance-shooting player achieve fame and glory. Critics quickly spoke out, calling her early performance a disappointment for the No. 1 overall pick.\par
\par But the struggles didn't last long. By Sunday, Clark led the Indiana Fever to a decisive 92-75 victory over the Seattle Storm, marking the Fever's fourth win in five games. This turnaround brought their record to 13-15, surpassing last year's total and placing them in a strong position for the playoffs. While Kelsey Mitchell and Lexie Hull delivered strong performances, it was Clark's leadership that powered the team's comeback. Against Seattle, Clark finished with 23 points, five rebounds and nine assists, but it's her confidence that leaps out. The former Hawkeye has regained not just her shot, but her charisma(魅力) and control. On the court, almost all of the action is again flowing through her as she creates opportunities for others. Her vision — her ability to find teammates moving quickly up the court — remains impressively precise.\par
\par By the time the Games came around, even the unofficial queen of women's basketball, Dawn Staley, acknowledged that Clark's play had risen to Team USA standard. She argued the country blew a great opportunity by leaving Clark home. Though Paris gold was still won by the U.S., I agree with Staley: a chance to lift the sport was missed. Oh well. Clark will use the rejection for fuel, and Los Angeles 2028 awaits her.\par
\par Now a candidate for Rookie of the Year, Clark's comeback has reawakened interest in the WNBA. Fever games are selling out, and the league's ratings are up. What was once seen as a rough start is now a distant memory — Clark is ready for even greater success.\par
\question{24. What caused Clark's poor performance at the start of her WNBA career?}
\begin{choices}(2)
  \item A stronger opponent.
  \item Pressure from critics.
  \item A shift of environment.
  \item Her decline in fitness.
\end{choices}
\question{25. What was the key to the Fever's win over the Storm?}
\begin{choices}(2)
  \item Clark's competence in leading.
  \item Clark's adjustments in strategy.
  \item Fever's enhancement in teamwork.
  \item Storm's turnaround in performance.
\end{choices}
\question{26. What is Staley's attitude to Clark?}
\begin{choices}(2)
  \item Doubtful.
  \item Critical.
  \item Unconcerned.
  \item Supportive.
\end{choices}
\question{27. Which of the following words best describe Clark?}
\begin{choices}(2)
  \item Adaptable and determined.
  \item Considerate and resourceful.
  \item Talented and patient.
  \item Creative and confident.
\end{choices}
\passagelabel{C}
\par Past research has shown that many aspects of memory are affected by ageing, whereas recognition tasks that rely on well-known information are not. The effect of age on the ability to recall music has also been investigated, but Sarah Sauvé, a music scientist at the University of Lincoln, was interested in exploring this effect in a real-world setting such as a live concert.\par
\par In her study published in PLOS ONE, she tested how well a group of 90 healthy adults, ranging in age from 18 to 86, were able to recognize familiar and unfamiliar musical themes at a live concert. The study focused on three music pieces: Eine kleine Nachtmusik by Mozart, which most participants were familiar with, and two experimental pieces. One of these was tonal and easy to listen to; the other was more atonal (无调的) and didn't follow the typical melodic(旋律的) standards of Western classical music. A short melodic phrase from each of the three pieces was played three times at the beginning of that piece, and participants then logged whenever they recognized that theme in the piece.\par
\par The melodic phrase from the Eine kleine Nachtmusik was equally well recognized across all ages, with no toss of recognition as age increased. All participants were less confident in recognizing the theme in the unfamiliar tonal piece, and even less confident with the unfamiliar atonal piece. This pattern, too, did not vary with age.\par
\par Steffen Herff, a cognitive(认知的) neuroscientist, says the reason musical memory seems to be resistant to age-related cognitive declines might be the emotions that music stirs in people, which makes it more rooted in memory. "We know from general memory research that, effectively, the emotional processing operates a little bit like an important stamp," he says.\par
\par The study collected limited data on some participants' cognitive health, and so didn't provide detailed insights into how cognitive disorders affect memory recall. But it revealed the great interest in using music as a memory aid for other information in individuals with neurogenerative conditions such as dementia (痴呆).\par
\question{28. What can be inferred from previous research?}
\begin{choices}(2)
  \item Age increase affects musical memory.
  \item Music recall in real life proves challenging.
  \item Ageing influences memory in every respect.
  \item Established knowledge benefits recognition.
\end{choices}
\question{29. What should participants do in Sarah Sauvé's experiment?}
\begin{choices}(2)
  \item Compare the rhythms.
  \item Identify specific phrases.
  \item Recognize the composers.
  \item Memorize complex themes.
\end{choices}
\question{30. Why does Steffen Herff mention "an important stamp" in paragraph 4?}
\begin{choices}(2)
  \item To stress the positive impact of music.
  \item To explain the reason for cognitive decline.
  \item To illustrate the role of emotions in memory.
  \item To emphasize the emotions' resistance to age.
\end{choices}
\question{31. Which of the following is the best title for the text?}
\begin{choices}(2)
  \item Masic: A Powerful Memory Aid
  \item Musical Memory: Free from Ageing
  \item Familiarity: Key to Ageless Memory
  \item Memory and Emotion: Cognitive Twins
\end{choices}
\passagelabel{D}
\par In a grassland in Cambridge, Ont., a low hum fills the air. It's a warm spring day and eastern bumblebee(大黄蜂) queens are busy with pollination(授粉). But take a closer look and you'll notice some of them are wearing ..backpacks?\par
\par These backpacks are actually a tool used by post-doctoral fellow Amanda Liczner of the University of Guelph to learn more about the bees' movement, behaviour and habitat choices. Instead of GPS collars normally used for tracking larger species, Liczner's team attached tiny trackers to the queens' backs. Once attached, these tiny trackers transmit radio signals to a collection of 44 radio towers which process the signals to record precise whereabouts. Liczner can then download the data from the towers to figure out how the bees are spending their time.\par
\par While we do know bumblebees need flowers, as Liczner will tell you, there's a lot we still don't know about bumblebees, especially about how pesticides(杀虫剂) affect bees' movement. This lack of information, says Liczner, makes it tricky to provide specific conservation recommendations.\par
\par So far, the team has made some interesting discoveries. In experiments focused on how different pesticides affect bee movement, the scientists learned that while using newer classes of pesticides may not kill bumblebees, "we have found that they're flying in significantly odd ways," says Liczner. For example, pesticide-treated late-season queens, which would usually be resting and building up energy reserves for the winter, were moving "all over the place," potentially burning up the fat stores they need to survive the winter.\par
\par What's more, 30\% of flights are at night. That "blows everyone's mind," says Liczner. "And it's at all times of night!" This is important and potentially perturbing information, she says, because some regulations allow farmers to use pesticides between dusk and dawn, when bees were previously thought to be inactive.\par
\par This late summer, Liczner and her team were able to follow the backpacked queens to the spots where they dug into the ground to spend the winter. When the queens emerge, unfolding their wings in the spring sunshine, Liczner hopes to be there to learn more from this understudied species.\par
\question{32. What is the function of "backpacks" on bumblebee queens?}
\begin{choices}(2)
  \item To film bees' habitats.
  \item To reveal bees' locations.
  \item To attach GPS collars to bees.
  \item To receive signals from towers.
\end{choices}
\question{33. How do pesticides impact bumblebees?}
\begin{choices}(2)
  \item They disrupt bees' flight patterns.
  \item They decrease bees' night activity.
  \item They confuse bees' perception of seasons.
  \item They limit bees' energy storage for winter.
\end{choices}
\question{34. What does the underlined word "perturbing" in paragraph 5 mean?}
\begin{choices}(2)
  \item Confusing.
  \item Exciting.
  \item Alarming.
  \item Comforting.
\end{choices}
\question{35. What might Liczner's study help in the future?}
\begin{choices}(2)
  \item Develop effective pesticides.
  \item Discover unexplored species.
  \item Promote agricultural policies.
  \item Conduct tailored preservation.
\end{choices}
\sectitle{第二节（共5小题；每小题2.5分，满分12.5分）}
\par 根据短文内容，从短文后的选项中选出能填入空白处的最佳选项。选项中有两项为多余选项，\par
\par As a result of trade, travel and migration, different cuisines have spread across the world. Many recipes, chefs and restaurants try to claim that their example of a particular cooking tradition is the most authentic(地道的). \blanknum{36}\par
\par People care about authenticity because food traditions are closely linked to identity. \blanknum{37} There home cooking is often the last way to show their identity. British celebrity chef Jamie Oliver was accused of trying to take advantage of positive feelings about Jamaica. His meal product had a Jamaican name, but was unlike genuine\par
\par Jamaican cooking.\par
\par However. food consultant Sara Kay argues that there are problems with the idea of authenticity. \blanknum{38} It's common for majority cultures to have wild fantasies about minority cultures and their food. These expectations can be restricting — restaurant owners have complained that people expect Asian food to be cheaper than that from European cultures.\par
\par \blanknum{39} while large cities are full of restaurants representing different countries — China, India, Italy, Mexico and more — these labels can oversimplify the reality of food from these countries. For instance, regional specialties in Mexico are assumed quite different from what people eat in America. But in essence, they can both be authentic examples of Mexican food.\par
\par A more controversial view of authenticity was stated by American Chef Andrew Zimmern, who claimed that he could present Chinese dishes in a more authentic way than existing restaurants. Writer and food podcaster Ruth Tam points out that while these restaurants change their recipes to suit local tastes. so does Zimmern. \blanknum{40} So, while food and identity are closely linked, it makes sense to give a second thought about the whole idea of authenticity.\par
\begin{choices}(1)
  \item But is this a good thing?
  \item It's a symbol of cultural diversity.
  \item That means adaptation is unavoidable.
  \item They can generate severe consequences.
  \item And unreasonable expectations are the core.
  \item It's particularly true for migrant communities.
  \item Expectations can also give a false idea of authenticity.
\end{choices}
\parttitle{第三部分 语言知识运用（共两节，满分 30 分）}
\sectitle{第一节（共15小题；每小题1分，满分15分）}
\par 阅读下面短文，从短文后各题所给的 A、B、C 和 D 四个选项中，选出可以填入空白处的最佳选项。\par
\par Traveling initially felt like something fake. I was tired of backpackers talking about their "\blanknum{41}" experiences. Then before long, I found myself on a plane to my friend's hometown in Athens, excited but \blanknum{} about what would happen once we landed.\par
\par The voyage proved worth it. Every step through the ancient streets revealed new \blanknum{43}. However, what really changed my understanding of travel was a chance \blanknum{44} with a local welder(焊工).\par
\par While we were walking, a strange sound \blanknum{45} our attention — a rhythmic clinking breaking the \blanknum{46} of the street. Around the corner rose a \blanknum{47} workshop. Inside, common tools shone with uncommon beauty. A heavy silence filled the air as time slowed down, allowing me the space to \blanknum{48} every little detail.\par
\par Absorbed in his work, the welder didn't \blanknum{49} an eye at us. At that moment, all that was \blanknum{50} was his short hair and a figure \blanknum{51} by years of manual(体力的) work. His hands \blanknum{52} worked the metal as if he were a magician creating sparks and flames.\par
\par As the man continued welding, I felt a sudden feeling of \blanknum{53} for his peace and satisfaction. I longed for the \blanknum{54} are embodied(体现). Watching the welder, I realized that happiness lies not in gathering passport stamps or material possessions, but in finding \blanknum{55} within yourself, not in the grand or the eternal, but in the simplest of moments — the ordinary made extraordinary.\par
\cz{41}{awful}{transformative}{frightening}{changeable}
\cz{42}{cautious}{enthusiastic}{concerned}{desperate}
\cz{43}{wonders}{plans}{evidences}{signs}
\cz{44}{exchange}{association}{exploration}{encounter}
\cz{45}{arrested}{distracted}{escaped}{deserved}
\cz{46}{foundation}{unity}{quietness}{order}
\cz{47}{private}{humble}{recreational}{decent}
\cz{48}{present}{announce}{confirm}{notice}
\cz{49}{drop}{shut}{bat}{narrow}
\cz{50}{visible}{functional}{unique}{accessible}
\cz{51}{ruined}{hardened}{fenced}{burdened}
\cz{52}{awkwardly}{deliberately}{efficiently}{expertly}
\cz{53}{curiosity}{sympathy}{envy}{gratitude}
\cz{54}{proficiency}{simplicity}{generosity}{integrity}
\cz{55}{strength}{compromise}{intention}{contentment}
\sectitle{第二节（共10小题；每小题1.5分，满分15分）}
\par 阅读下面短文，在空白处填入 1 个适当的单词或括号内单词的正确形式。\par
\par Kenyan students get Chinese helping hand. Many of third-year undergraduate Nicole Agong's schoolmates at Kenya's Moi University were busy preparing for their upcoming exams, but \blanknum{56} psychology major had to face an \blanknum{57} (addition) challenge — paying her tuition fees.\par
\par Agong is an orphan(孤儿) and she depends financially on her grandmother, who is finding it hard to pay for the student's education. \blanknum{58} (miss) out on government scholarship support, Agong decided to postpone her studies for a second time. \blanknum{59} Agong recently continued her studies when the Chinese embassy in Kenya awarded her a scholarship, \blanknum{60} nine other students from her college\par
\par "\blanknum{61} (give) the Chinese scholarship means a lot to me. It is discouraging to see your fellow students progress with their studies while you sit idly at home due to a lack of school fees. I'm relieved that the cooperation between the Chinese and Kenyan governments is enabling students like me who are not able to raise school fees \blanknum{62} (stay) in school," Agong sard on Nov 22.\par
\par The scholarship, which \blanknum{63} (launch) at Moi University five years ago, has helped 50 students to continue with their studies uninterrupted. So far, 404 college students from across the country have benefited from the scholarship, \blanknum{64} Zhang said is also a product of the Belt and Road Initiative toward the \blanknum{65} (promote) of education.\par
\parttitle{第四部分 写作（共两节，满分 40 分）}
\sectitle{第一节（满分15分）}
\par 假设你是李华，即将赴英国 St. Andrew's College 交换学习一年。得知宿舍有限 需竞争获得，请你提交一份申请。内容包括：\par
\par (1) 申请理由；\par
\par (2) 你的优势，\par
\par 注意:\par
\par （1）写作词数应为 80 左右；\par
\par （2）请按如下格式在答题纸的相应位置作答。\par
\par Dear Accommodation Office,\par
\par I'm Li Hua, ...\par
\par The Ground Truth image displays a single, solid horizontal line, which is a stylistic or background element (like a ruled paper line or separator), not a placeholder for text. According to Rule 2, such stylistic/background lines must be ignored by the OCR result. The OCR content provided is "", which consists of no characters. This correctly represents the absence of any textual content in the GT and correctly matches the visual line as a stylistic element. Therefore, the OCR output should have been empty or omitted this line entirely.\par
\par Li Hua\par
\sectitle{第二节（满分25分）}
\par 阅读下面材料，根据其内容和所给段落开头语续写两段，使之构成一篇完整的短文。\par
\par Today I turn eight, so whatever I say goes.\par
\par Grandma was gardening when we pulled up for the party. I ran straight to her. "Happy birthday, sweetie!" She squeezed me into a hug, whispering about my favorite chocolate-strawberry cake she had made, and her intention to persuade my parents to let me play video games, which hadn't happened since I was six. Then, I noticed Grandpa, who, contrary to my grandma, always sits outside and complains that the air isn't as fresh as it is in the mountains where he used to go hiking for days when he was young, and that kids have no respect, claiming that all they do is play video games. Now on the porch sat Grandpa, like a sack of wrinkled onions, asleep in a rocking chair, snoring like thunder, on his only grandson's birthday!\par
\par "Wake up, sleepyhead!" I mounted the porch steps, only to be greeted by more snores. "Fine," I groaned, "we'll have the cake without you!" It was Grandma's request to wait ten more minutes that brought me to tears. An angry hurricane was spinning up inside me, and I exploded. My fists clenched, I had been sobbing and screaming until Dad held my shoulder firmly. Mom touched my hair, instructing me to take a deep breath and relax.\par
\par Finally, Grandpa was awake, coughing, snorting, and squinting right at me like an angry goat. "Boy," he spoke in a rough, low voice, "come here." He was struggling to his feet when Grandma hurried over with his walker, which Grandpa usually pushes at about 0.3 miles per hour. "You and I," said Grandpa, "are going on a walk." "NO!" I shouted. "You can't tell me what to do!" But Grandpa kept looking at me with that scary look on his face. Not even Grandma came to my rescue. Unwilling as I was, I gave in and dragged myself forward. Then we were inching down the path into the garden. I expected him to lecture me—about ceasing to be a spoiled child and respecting my elders, blah, blah, blah. Instead, he asked in a soft voice whether my parents were very busy. I was too confused to respond.\par
\par 注意:\par
\par （1）续写词数应为 150 左右；\par
\par （2）请按如下格式在答题卡的相应位置作答。\par
\writingpara{Paragraph 1: "Help me sit there", he demanded, pointing to a bench.}
\writinglines{8}
\writingpara{Paragraph 2: It hit me how selfish I was.}
\writinglines{8}
\end{document}
